%%%%%%%%%%%%%%%%%%%%%%%%%%%%%%%%%%%%%%%%%%%%%%%%%%%%%%%%%%%%%%%%%%%%%%%%%%%%%%%
% APPENDIX BB: FORMAL-INTERVENTION-EQUILIBRIA
% Operationalizing Counterfactual Equilibria & Tipping Points
%%%%%%%%%%%%%%%%%%%%%%%%%%%%%%%%%%%%%%%%%%%%%%%%%%%%%%%%%%%%%%%%%%%%%%%%%%%%%%%

\documentclass[11pt, a4paper]{article}
\usepackage[utf8]{inputenc}
\usepackage[margin=1in]{geometry}
\usepackage{amsmath, amssymb, amsthm}
\usepackage{booktabs, array, multirow}
\usepackage{xcolor}
\usepackage{tcolorbox}
\usepackage{enumitem}
\usepackage{hyperref}

\title{Appendix BB: Formal Framework for Intervention Equilibria and Tipping Points}
\author{Evidence-Based Framework for Economic and Social Behavior}
\date{January 2026}

\begin{document}

\maketitle

%------------------------------------------------------------------------------
% METADATA & ABSTRACT
%------------------------------------------------------------------------------

\begin{tcolorbox}[colback=gray!10!white, colframe=gray!75!black, title=Appendix Metadata]
\textbf{Code:} BB \\
\textbf{Category:} FORMAL-INTERVENTION-EQUILIBRIA \\
\textbf{Version:} 1.0 (January 2026) \\
\textbf{Purpose:} Formalize the operationalization of equilibrium stability $\xi$ and critical threshold $\kappa_{\text{crit}}$
for practitioner-level intervention design \\
\textbf{Primary Axioms:} W11-W13 (Dynamic Modulation, Boundary Conditions, Probability Bridge) \\
\textbf{Cross-References:} Ch14 (BCS), Ch15 (Integration), AV (Willingness),
AW (BCJ), BA (Segmentation), \textbf{XVIII (LIT-FEHR-METHOD)} \\
\end{tcolorbox}

\begin{abstract}
This appendix provides the formal framework for operationalizing counterfactual behavioral change equilibria.
Starting from the Willingness-Function (W11-W13), we establish two key diagnostic parameters:
(1) equilibrium stability $\xi$ measuring how inherently robust the status quo is, and
(2) critical threshold $\kappa_{\text{crit}}$ defining where the system enters a tipping-point regime.
We develop a practical diagnostic questionnaire allowing practitioners to estimate these parameters,
then map them to optimal intervention strategies. The framework synthesizes theory (formal mathematics),
measurement (practitioner tools), and application (worked examples) for the Behavioral Change Integration model (Chapter 15).
\end{abstract}

%------------------------------------------------------------------------------
% FUNDAMENTAL QUESTION
%------------------------------------------------------------------------------

\begin{tcolorbox}[colback=blue!5!white, colframe=blue!75!black,
    title=\textbf{Fundamental Question}]
\textbf{Given a behavioral change context, how do we operationalize equilibrium stability $\xi$
and critical threshold $\kappa_{\text{crit}}$ to determine optimal intervention timing and intensity?}

\medskip
\textit{This appendix answers: What mathematical framework enables practitioners to diagnose
when a behavioral system is at risk of tipping and what intervention strategies are optimal
for each equilibrium regime?}
\end{tcolorbox}

%------------------------------------------------------------------------------
% METHODOLOGICAL NOTE: EXCLUSION PRINCIPLE (XVIII)
%------------------------------------------------------------------------------

\begin{tcolorbox}[colback=red!5!white, colframe=red!75!black,
    title=\textbf{Methodological Note: Exclusion Principle (Appendix XVIII)}]

\textbf{BB and the Null Hypothesis:}

This appendix formalizes equilibrium dynamics that \textit{only emerge when complementarity is present}.
Following Appendix XVIII (LIT-FEHR-METHOD), we note:

\begin{itemize}[nosep]
    \item \textbf{If $\gamma = 0$:} All $\kappa$-modulations ($\kappa_{AWX}, \kappa_{KON}, \kappa_{JNY}$) collapse to additive effects. The non-linear tipping dynamics disappear---system becomes linear.
    \item \textbf{If $\gamma \neq 0$:} Complementarity creates the non-linear $\lambda(\kappa)$ dynamics that generate true tipping points ($\kappa_{\text{crit}}$).
\end{itemize}

\medskip
\textbf{Implication for Practitioners:}
\begin{itemize}[nosep]
    \item The hyperbolic singularity (Q4 regime) exists \textit{only because} interventions interact non-additively
    \item Before applying BB's tipping-point framework, verify that $\gamma \neq 0$ holds in your context
    \item If complementarity cannot be established, simpler linear intervention models may suffice
\end{itemize}

\textit{BB operationalizes dynamics that are conditional on the rejection of the additive null hypothesis.}
\end{tcolorbox}

%------------------------------------------------------------------------------
% TABLE OF CONTENTS
%------------------------------------------------------------------------------

\section*{Quick Reference: Appendix Structure}

\begin{itemize}
    \item \textbf{Section 1:} Three Counterfactual Equilibria (Baseline, Drifting SQ, Goal)
    \item \textbf{Section 2:} Equilibrium Stability $\xi$ — Operationalization Framework
    \item \textbf{Section 3:} Critical Threshold $\kappa_{\text{crit}}$ — Domain-Dependent Definition
    \item \textbf{Section 4:} The Five Mechanisms Determining $\kappa_{\text{crit}}$
    \item \textbf{Section 5:} Mathematical Model: Non-Linear Decay with Context-Dependent $\lambda(\kappa)$
    \item \textbf{Section 6:} Diagnostic Tool for Practitioners (Questionnaire + Calculator)
    \item \textbf{Section 7:} Decision Matrix: $(\xi, \kappa_{\text{crit}}) \to$ Intervention Strategy
    \item \textbf{Section 8:} Worked Examples (Solar, COVID, Smoking, Workplace)
    \item \textbf{Section 9:} Integration with Chapter 15 (Intervention Design Workflow)
\end{itemize}

%=============================================================================
% SECTION 1: THREE COUNTERFACTUAL EQUILIBRIA
%=============================================================================

\section{The Three Counterfactual Equilibria}

\subsection{Definition: What We Compare}

An \textbf{equilibrium} in the context of behavioral change is a trajectory $\kappa(t)$
describing how the effective behavioral willingness evolves over time under a specific
intervention regime $I(t)$.

\begin{definition}[Three Counterfactual Equilibria]
For any behavioral change context, we define three equilibria:

\begin{enumerate}
    \item \textbf{Baseline Equilibrium} ($I = 0$):
    No intervention, natural deterioration to collapse
    $$\frac{d\kappa}{dt} = -\lambda(\kappa) \cdot \kappa(t), \quad \kappa(0) = \kappa_{\text{SQ}}$$
    Outcome: $\kappa(t) \to 0$ exponentially or hyperbolically

    \item \textbf{Status Quo Equilibrium (Drifting)} ($I = I_{\text{min}} \cdot 0.8$):
    Minimal intervention insufficient to halt drift; gradual deterioration
    $$\frac{d\kappa}{dt} = -\lambda(\kappa) \cdot \kappa(t) + I_{\text{drift}}, \quad \kappa(0) = \kappa_{\text{SQ}}$$
    where $I_{\text{drift}}$ is 80\% of minimum required
    Outcome: $\kappa(t)$ decays slowly toward $\kappa_{\text{crit}}$ over 12-18 months

    \item \textbf{Stationary Status Quo} ($I = I_{\text{min}}$):
    Minimum intervention exactly balances decay; stability maintained
    $$\frac{d\kappa}{dt} = 0, \quad \kappa(t) = \kappa_{\text{SQ}} \, \forall t$$
    where $I_{\text{min}} = \lambda(\kappa_{\text{SQ}}) \cdot \kappa_{\text{SQ}}$ (by definition)
    Outcome: $\kappa(t)$ remains constant forever

    \item \textbf{Goal Equilibrium} ($I = I_{\text{opt}}$):
    Intensive intervention drives improvement toward target
    $$\frac{d\kappa}{dt} = +\mu \cdot (\kappa_{\text{max}} - \kappa(t)), \quad \kappa(0) = \kappa_{\text{SQ}}$$
    Outcome: $\kappa(t)$ grows to $\kappa_{\text{goal}}$ with characteristic time $\tau_{\text{ramp}}$
\end{enumerate}
\end{definition}

\subsection{Visualization: The Four Regimes}

\begin{equation}
\text{Table 1: Four Equilibrium Regimes}
\end{equation}

\begin{center}
\begin{tabular}{|l|c|c|c|c|}
\hline
\textbf{Regime} & \textbf{Intervention} & \textbf{Dynamics} & \textbf{Outcome} & \textbf{Timescale} \\
\hline
Baseline & $I = 0$ & Exponential/Hyperbolic decay & $\kappa \to 0$ & 2-12 mo \\
Drifting SQ & $I = 0.8 I_{\min}$ & Slow decay + drift & $\kappa \to \kappa_{\text{crit}}$ & 12-18 mo \\
Stationary SQ & $I = I_{\min}$ & $d\kappa/dt = 0$ & $\kappa = \text{const}$ & $\infty$ \\
Goal & $I = I_{\text{opt}}$ & Logistic growth & $\kappa \to \kappa_{\text{goal}}$ & 3-12 mo \\
\hline
\end{tabular}
\end{center}

%=============================================================================
% SECTION 2: EQUILIBRIUM STABILITY ξ
%=============================================================================

\section{Operationalizing Equilibrium Stability: $\xi$}

\subsection{Definition}

\begin{definition}[Equilibrium Stability $\xi$]
Equilibrium stability $\xi \in [0, 1]$ measures how inherently robust the current
status quo is to deterioration, independent of active intervention.

$$\xi = f(\text{Network Effects}, \text{Institutional Support}, \text{Path Dependence})$$

\vspace{0.5cm}
\noindent
where each component is normalized to $[0, 1]$ scale:

\begin{align}
\xi &= 0 \quad \text{(fully labile: system collapses without any support)} \\
\xi &= 0.5 \quad \text{(metastable: ambivalent, could go either way)} \\
\xi &= 1 \quad \text{(fully stable: self-reinforcing, needs little support)}
\end{align}
\end{definition}

\subsection{Three Components of $\xi$}

\subsubsection{Component 1: Network Effects Strength ($N_{\text{net}}$)}

\begin{equation}
N_{\text{net}} = \text{Proportion of peers who adopted} \times \text{Influence per peer}
\end{equation}

\begin{itemize}
    \item \textbf{High} ($N_{\text{net}} \geq 0.5$):
    Neighbors/colleagues visibly adopt; social proof is strong
    \\ Example: Solar adoption in neighborhood with 40\% penetration

    \item \textbf{Medium} ($0.2 \leq N_{\text{net}} < 0.5$):
    Some peer influence, but not dominant
    \\ Example: COVID vaccination in mixed community (60\% vax, 40\% hesitant)

    \item \textbf{Low} ($N_{\text{net}} < 0.2$):
    Peer adoption invisible or irrelevant
    \\ Example: Zahnseide brushing (no one sees if you floss)
\end{itemize}

\subsubsection{Component 2: Institutional Support ($I_{\text{struct}}$)}

\begin{equation}
I_{\text{struct}} = \text{Policy} + \text{Infrastructure} + \text{Social Norm}
\end{equation}

\begin{itemize}
    \item \textbf{High} ($I_{\text{struct}} \geq 0.5$):
    Policy, infrastructure, and norms all support behavior
    \\ Example: Solar adoption with subsidies, grid connection, public support

    \item \textbf{Medium} ($0.2 \leq I_{\text{struct}} < 0.5$):
    Some structural support
    \\ Example: EV adoption with charging infrastructure in some cities

    \item \textbf{Low} ($I_{\text{struct}} < 0.2$):
    Little structural support; individual must carry burden
    \\ Example: Smoking cessation (individual responsibility)
\end{itemize}

\subsubsection{Component 3: Path Dependence ($P_{\text{dep}}$)}

\begin{equation}
P_{\text{dep}} = \text{How long has behavior been practiced?} \times \text{Integration into identity/routine}
\end{equation}

\begin{itemize}
    \item \textbf{High} ($P_{\text{dep}} \geq 0.5$):
    Behavior is deeply embedded in routine, identity, or environmental design
    \\ Example: Long-term smoker, established solar adopter

    \item \textbf{Medium} ($0.2 \leq P_{\text{dep}} < 0.5$):
    Behavior is moderately established
    \\ Example: 2-year gym member, new parent with established routines

    \item \textbf{Low} ($P_{\text{dep}} < 0.2$):
    Behavior is recent or surface-level
    \\ Example: One-time COVID vaccination decision
\end{itemize}

\subsection{Aggregate $\xi$ Calculation}

\begin{equation}
\xi = \frac{1}{3}(N_{\text{net}} + I_{\text{struct}} + P_{\text{dep}}) + \epsilon_{\text{interaction}}
\end{equation}

\noindent
where $\epsilon_{\text{interaction}}$ captures synergies (e.g., strong networks
amplify institutional support):

\begin{equation}
\epsilon_{\text{interaction}} = 0.1 \times N_{\text{net}} \times I_{\text{struct}}
\end{equation}

\subsection{Interpretation}

\begin{center}
\begin{tabular}{|c|c|l|}
\hline
$\xi$ Range & Label & Interpretation \\
\hline
$0.0 - 0.3$ & Labile & System highly fragile; will collapse without intervention \\
$0.3 - 0.6$ & Metastable & System ambivalent; could go either direction \\
$0.6 - 0.8$ & Stable & System self-reinforcing; needs moderate support \\
$0.8 - 1.0$ & Highly Stable & System is robust; minimal support needed \\
\hline
\end{tabular}
\end{center}

%=============================================================================
% SECTION 3: CRITICAL THRESHOLD κ_crit
%=============================================================================

\section{Operationalizing Critical Threshold: $\kappa_{\text{crit}}$}

\subsection{Definition}

\begin{definition}[Critical Threshold $\kappa_{\text{crit}}$]
The critical threshold $\kappa_{\text{crit}}$ is the value below which:

\begin{enumerate}
    \item The system enters a regime of accelerating decay (bifurcation point)
    \item Dropout/relapse probability exceeds 50\% within domain-specific horizon
    \item Intervention becomes increasingly ineffective (marginal returns diminish)
\end{enumerate}

Mathematically:
\begin{equation}
\kappa_{\text{crit}} = \arg\min_\kappa \left\{ \frac{d^2\kappa}{dt^2} \text{ becomes sharply negative} \right\}
\end{equation}

Or operationally:
\begin{equation}
\kappa_{\text{crit}} = \inf\{\kappa : P(\text{abandon} | \kappa(t), I=I_{\min}) > 0.5 \text{ within } \tau_{\text{horizon}}\}
\end{equation}
\end{definition}

\subsection{Domain-Specific Values (Empirical Reference)}

\begin{center}
\begin{tabular}{|l|c|c|c|}
\hline
\textbf{Domain} & $\kappa_{\text{crit}}$ & \textbf{$\tau_{\text{horizon}}$} & \textbf{Confidence} \\
\hline
Solar Adoption & 0.22 & 6 months & High (50+ studies) \\
COVID Vaccination & 0.10 & 3 months & Low (highly metastable) \\
Smoking Cessation & 0.38 & 2 weeks & High (100+ studies) \\
Workplace Wellness & 0.28 & 3 months & Medium (20+ studies) \\
Dental Flossing & 0.32 & 2 weeks & Low (limited data) \\
EV Adoption (with infrastructure) & 0.20 & 6 months & Medium \\
EV Adoption (without infrastructure) & 0.45 & 1 month & Medium \\
\hline
\end{tabular}
\end{center}

\subsection{The Five Mechanisms Determining $\kappa_{\text{crit}}$}

The actual value of $\kappa_{\text{crit}}$ for a given domain emerges from five behavioral
mechanisms working in concert. Each mechanism weights how fragile the system is:

\subsubsection{Mechanism M1: Social Network Structure ($w_1$)}

\begin{equation}
w_1 = N_{\text{net}} \times \text{Cascading Effect}
\end{equation}

\begin{itemize}
    \item \textbf{High cascade effect} (M1 = strong):
    Each adopter recruits more adopters exponentially
    \\ $\kappa_{\text{crit}}$ decreases (tipping point comes earlier)
    \\ Example: Solar adoption, COVID vaccination

    \item \textbf{Low cascade effect} (M1 = weak):
    Adoption is independent per person
    \\ $\kappa_{\text{crit}}$ increases (tipping point comes later)
    \\ Example: Flossing, personal exercise habits
\end{itemize}

\textbf{Effect on $\kappa_{\text{crit}}$:} Increase in $w_1 \Rightarrow$ decrease in $\kappa_{\text{crit}}$

\subsubsection{Mechanism M2: Habit Formation Depth ($w_2$)}

\begin{equation}
w_2 = \text{Years of daily practice} \times \text{Neurological Entrenchment}
\end{equation}

\begin{itemize}
    \item \textbf{Deep habits} (M2 = strong):
    Decades of daily practice; extinction burst when breaking
    \\ $\kappa_{\text{crit}}$ increases (system is superficially stable but brittle)
    \\ Example: Smoking (sharp relapse), Dental flossing (forgotten slowly then forgotten)

    \item \textbf{Shallow habits} (M2 = weak):
    Recent or infrequent; fades gradually
    \\ $\kappa_{\text{crit}}$ decreases (but decay is slow)
    \\ Example: COVID vaccination (one-time event), Solar adoption (one-time decision)
\end{itemize}

\textbf{Effect on $\kappa_{\text{crit}}$:} Increase in $w_2 \Rightarrow$ increase in $\kappa_{\text{crit}}$
(but with sharp bifurcation)

\subsubsection{Mechanism M3: Environmental Feedback ($w_3$)}

\begin{equation}
w_3 = \text{Clarity of outcome} \times \text{Immediacy of feedback}
\end{equation}

\begin{itemize}
    \item \textbf{Strong feedback} (M3 = strong):
    Clear, immediate signals that behavior works
    \\ $\kappa_{\text{crit}}$ decreases (positive reinforcement holds $\kappa$ up)
    \\ Example: Solar (see electricity bill drop immediately), Fitness (feel improvement in weeks)

    \item \textbf{Weak feedback} (M3 = weak):
    Delayed or ambiguous signals
    \\ $\kappa_{\text{crit}}$ increases (nothing keeps $\kappa$ up during waiting period)
    \\ Example: COVID vax (protection is unseen), Climate action (effects in decades)
\end{itemize}

\textbf{Effect on $\kappa_{\text{crit}}$:} Increase in $w_3 \Rightarrow$ decrease in $\kappa_{\text{crit}}$

\subsubsection{Mechanism M4: Cognitive Lock-in ($w_4$)}

\begin{equation}
w_4 = \text{Identity binding} \times \text{Consistency motivation}
\end{equation}

\begin{itemize}
    \item \textbf{Strong lock-in} (M4 = strong):
    Identity is bound to behavior; cognitive dissonance maintains it
    \\ $\kappa_{\text{crit}}$ decreases (identity keeps $\kappa$ up even if behavior wavers)
    \\ Example: COVID vax ("I am vaccinated" is identity), Environmentalist adopting solar

    \item \textbf{Weak lock-in} (M4 = weak):
    Behavior is instrumental, not identity-defining
    \\ $\kappa_{\text{crit}}$ increases (behavior fades when mood/motivation dips)
    \\ Example: Flossing (no one asks "are you a flosser?"), Generic gym-goer
\end{itemize}

\textbf{Effect on $\kappa_{\text{crit}}$:} Increase in $w_4 \Rightarrow$ decrease in $\kappa_{\text{crit}}$

\subsubsection{Mechanism M5: Institutional Support ($w_5$)}

\begin{equation}
w_5 = \text{Policy} + \text{Infrastructure} + \text{Social Norms}
\end{equation}

\begin{itemize}
    \item \textbf{Strong support} (M5 = strong):
    Institutions make behavior easy and default
    \\ $\kappa_{\text{crit}}$ decreases (structure compensates for individual wavering)
    \\ Example: Solar (grid connection is automatic, subsidy pays for it),
               Workplace (facilities provided, time allocated)

    \item \textbf{Weak support} (M5 = weak):
    Individual bears full burden
    \\ $\kappa_{\text{crit}}$ increases (behavior depends entirely on personal motivation)
    \\ Example: Smoking cessation (no infrastructure), Flossing (no social support)
\end{itemize}

\textbf{Effect on $\kappa_{\text{crit}}$:} Increase in $w_5 \Rightarrow$ decrease in $\kappa_{\text{crit}}$

\subsection{Aggregate Formula for $\kappa_{\text{crit}}$}

\begin{equation}
\kappa_{\text{crit}} = \kappa_0 - (w_1 + w_3 + w_4 + w_5) + w_2 + \epsilon_{\text{domain}}
\end{equation}

\noindent
where:
\begin{itemize}
    \item $\kappa_0 = 0.30$ is the baseline threshold
    \item $w_i \in [0, 0.15]$ are the normalized mechanism weights
    \item $\epsilon_{\text{domain}} \in [-0.05, 0.05]$ is domain-specific adjustment
\end{itemize}

\textbf{Interpretation:}
\begin{itemize}
    \item Social cascades ($w_1$), feedback ($w_3$), identity ($w_4$), and institutions ($w_5$)
    all \textit{lower} $\kappa_{\text{crit}}$ (tipping point comes earlier because they help maintain $\kappa$)

    \item Deep habits ($w_2$) \textit{raise} $\kappa_{\text{crit}}$ (tipping point is higher because
    withdrawal is more intense)
\end{itemize}

%=============================================================================
% SECTION 4: NON-LINEAR DECAY DYNAMICS
%=============================================================================

\section{Mathematical Model: Non-Linear Decay with Context-Dependent $\lambda(\kappa)$}

\subsection{The Fundamental Equation}

\begin{equation}
\frac{d\kappa}{dt} = -\lambda(\kappa) \cdot \kappa(t) + I(t)
\end{equation}

\noindent
where:
\begin{itemize}
    \item $\kappa(t) \in [0, 1]$ is effective behavioral willingness (from W11 modulation)
    \item $\lambda(\kappa) > 0$ is the context-dependent drift rate (non-linear!)
    \item $I(t) \geq 0$ is intervention intensity
\end{itemize}

\subsection{Context-Dependent Forms of $\lambda(\kappa)$}

The form of $\lambda(\kappa)$ depends on $(\xi, \text{Distance to Optimum})$:

\subsubsection{Regime Q1: Stable + Near Optimum ($\xi > 0.6$ \& $\delta > 0.6$)}

\begin{equation}
\lambda(\kappa) = \lambda_0^{(I)} \cdot \left[1 + \beta \cdot \frac{\kappa_{\max} - \kappa}{\kappa_{\max}}\right]
\end{equation}

\noindent
\textbf{Characteristics:}
\begin{itemize}
    \item Polynomial growth (smooth acceleration as $\kappa$ decays)
    \item No singularities; drift is bounded
    \item Tipping point exists but is "soft" (gradual transition)
\end{itemize}

\noindent
\textbf{Interpretation:} "Good equilibrium near optimum — small interventions work well"

\subsubsection{Regime Q2/Q3: Mixed Stability or Proximity}

\begin{equation}
\lambda(\kappa) = \lambda_0^{(II)} \cdot \exp\left[\gamma \cdot \frac{\kappa_{\max} - \kappa}{\kappa_{\max}}\right]
\end{equation}

\noindent
\textbf{Characteristics:}
\begin{itemize}
    \item Exponential growth (dramatic acceleration as $\kappa$ decays)
    \item Still no singularities but strong nonlinearity
    \item Tipping point is "hard" (rapid transition)
\end{itemize}

\noindent
\textbf{Interpretation:} "Mixed context — intervention needed before deterioration accelerates"

\subsubsection{Regime Q4: Labile + Far from Optimum ($\xi < 0.6$ \& $\delta < 0.6$)}

\begin{equation}
\lambda(\kappa) = \lambda_0^{(IV)} \cdot \frac{\kappa_{\text{crit}}}{\kappa - \kappa_{\text{crit}} + \epsilon}
\end{equation}

\noindent
\textbf{Characteristics:}
\begin{itemize}
    \item Hyperbolic singularity: $\lambda(\kappa) \to \infty$ as $\kappa \to \kappa_{\text{crit}}^+$
    \item True bifurcation/tipping point
    \item No smooth transition; system is unstable
\end{itemize}

\noindent
\textbf{Interpretation:} "Critical regime — decision point now; intervention or abandon"

%=============================================================================
% SECTION 5: DIAGNOSTIC TOOL FOR PRACTITIONERS
%=============================================================================

\section{Diagnostic Tool: Questionnaire + Calculator}

This section provides a practical tool for practitioners to estimate $(\xi, \kappa_{\text{crit}}, \lambda(\kappa))$
for their specific behavioral change context.

\subsection{Step 1: Measure Equilibrium Stability ($\xi$)}

Answer the following questions on a 1-5 scale (1=Low, 5=High):

\begin{enumerate}
    \item \textbf{Network Effects:}
    Among your target population, what percentage have already adopted the behavior?
    \begin{itemize}
        \item 1: $< 5\%$ adoption (no social proof yet)
        \item 3: 15-30\% adoption (some peer influence)
        \item 5: $> 40\%$ adoption (strong social proof)
    \end{itemize}
    \textit{Store as:} $Q_1 \in [1, 5]$

    \item \textbf{Cascading Effect:}
    When one person adopts, does it visibly influence others to adopt?
    \begin{itemize}
        \item 1: No cascades (adoption is independent)
        \item 3: Some cascades (maybe 10-20\% of neighbors follow)
        \item 5: Strong cascades (40\%+ of neighbors follow)
    \end{itemize}
    \textit{Store as:} $Q_2 \in [1, 5]$

    \item \textbf{Institutional Support:}
    Is there policy, infrastructure, or social norm support?
    \begin{itemize}
        \item 1: None (individual bears full burden)
        \item 3: Partial (some infrastructure or weak policy)
        \item 5: Strong (subsidies, infrastructure, norm support all present)
    \end{itemize}
    \textit{Store as:} $Q_3 \in [1, 5]$

    \item \textbf{Path Dependence:}
    How long has this behavior typically been practiced?
    \begin{itemize}
        \item 1: New ($< 3$ months)
        \item 3: Moderate (6-12 months)
        \item 5: Established ($> 2$ years, deeply embedded in routine)
    \end{itemize}
    \textit{Store as:} $Q_4 \in [1, 5]$
\end{enumerate}

\noindent
\textbf{Calculate $\xi$:}
\begin{equation}
\xi = \frac{1}{20} \left[ \frac{Q_1 + Q_2}{2} + Q_3 + Q_4 \right]
\end{equation}

\noindent
(Normalize to $[0, 1]$ scale)

\subsection{Step 2: Measure Mechanisms for $\kappa_{\text{crit}}$}

For each of the five mechanisms, rate on a 1-5 scale:

\begin{enumerate}
    \item \textbf{Social Network (M1):}
    Do people see/hear when others adopt? Is social proof visible?
    \begin{itemize}
        \item 1: No visibility (e.g., flossing)
        \item 3: Moderate visibility (e.g., gym membership)
        \item 5: High visibility (e.g., solar panels, electric cars)
    \end{itemize}
    \textit{Store as:} $M_1 \in [1, 5]$, Weight: $w_1 = -0.15 \times (M_1/5)$

    \item \textbf{Habit Depth (M2):}
    Is this a deep habit or a one-time decision?
    \begin{itemize}
        \item 1: One-time decision (e.g., vaccination)
        \item 3: Moderate habit (6-12 months practice)
        \item 5: Deep habit (20+ years daily practice)
    \end{itemize}
    \textit{Store as:} $M_2 \in [1, 5]$, Weight: $w_2 = +0.15 \times (M_2/5)$

    \item \textbf{Environmental Feedback (M3):}
    Is the success/benefit immediately visible?
    \begin{itemize}
        \item 1: No feedback or delayed (e.g., climate action)
        \item 3: Moderate feedback (e.g., fitness improvement in weeks)
        \item 5: Strong immediate feedback (e.g., solar bill drops monthly)
    \end{itemize}
    \textit{Store as:} $M_3 \in [1, 5]$, Weight: $w_3 = -0.15 \times (M_3/5)$

    \item \textbf{Cognitive Lock-in (M4):}
    Is this behavior part of someone's identity?
    \begin{itemize}
        \item 1: Not identity-defining (e.g., generic wellness)
        \item 3: Somewhat identity-defining (e.g., "environmentalist")
        \item 5: Core identity (e.g., "I am vaccinated" or "I am a smoker")
    \end{itemize}
    \textit{Store as:} $M_4 \in [1, 5]$, Weight: $w_4 = -0.15 \times (M_4/5)$

    \item \textbf{Institutional Support (M5):}
    Does the system make the behavior easy or default?
    \begin{itemize}
        \item 1: No structural support (individual must initiate everything)
        \item 3: Partial support (infrastructure in place, some friction remains)
        \item 5: Full support (behavior is default, friction removed)
    \end{itemize}
    \textit{Store as:} $M_5 \in [1, 5]$, Weight: $w_5 = -0.15 \times (M_5/5)$
\end{enumerate}

\noindent
\textbf{Calculate $\kappa_{\text{crit}}$:}
\begin{equation}
\kappa_{\text{crit}} = 0.30 + w_1 + w_2 + w_3 + w_4 + w_5
\end{equation}

\noindent
Clamp to $[0.05, 0.50]$ range (should be between these bounds).

\subsection{Step 3: Classify Your Context}

Based on $(\xi, \delta)$ where $\delta = $ (current adoption) / (optimal adoption),
determine which quadrant you are in:

\begin{center}
\begin{tabular}{|c|c|c|}
\hline
& $\delta > 0.6$ (Near Optimum) & $\delta \leq 0.6$ (Far from Optimum) \\
\hline
$\xi > 0.6$ (Stable) & \textbf{Q1:} MAINTAIN & \textbf{Q2:} PREVENT \\
$\xi \leq 0.6$ (Labile) & \textbf{Q3:} DEFEND & \textbf{Q4:} DECIDE \\
\hline
\end{tabular}
\end{center}

\subsection{Step 4: Determine $\lambda(\kappa)$ Form and Intervention Need}

Based on your quadrant:

\begin{itemize}
    \item \textbf{Q1 (Maintain):} $\lambda(\kappa) = $ POLYNOMIAL; small interventions wirken gut
    \\ \textit{Action:} Monthly reminders, annual events. Cost: Low.

    \item \textbf{Q2 (Prevent):} $\lambda(\kappa) = $ EXPONENTIAL; invest now before acceleration
    \\ \textit{Action:} Monthly + quarterly engagement. Cost: Medium.
    \\ \textit{Timeline:} Act NOW; window closes in 6 months.

    \item \textbf{Q3 (Defend):} $\lambda(\kappa) = $ EXPONENTIAL; proportional effort needed
    \\ \textit{Action:} Regular touchpoints (monthly or weekly). Cost: Medium-High.
    \\ \textit{Timeline:} Continuous engagement.

    \item \textbf{Q4 (Decide):} $\lambda(\kappa) = $ HYPERBOLIC; tipping point is near
    \\ \textit{Action:} Intensive intervention (weekly) OR accept loss. Cost: Very High OR Zero.
    \\ \textit{Timeline:} Decision must be made within 1-3 months.
\end{itemize}

%=============================================================================
% SECTION 6: DECISION MATRIX
%=============================================================================

\section{Decision Matrix: From Diagnosis to Intervention Strategy}

\begin{center}
\small
\begin{tabular}{|p{1.8cm}|p{1.8cm}|p{1.8cm}|p{1.8cm}|p{1.8cm}|}
\hline
\textbf{Quadrant} & \textbf{$\xi$, $\delta$} & \textbf{$\lambda(\kappa)$ Form} &
\textbf{Intervention Strategy} & \textbf{Estimated Cost} \\
\hline

Q1: MAINTAIN
& $\xi > 0.6$ $\delta > 0.6$
& Polynomial
& Monthly email + annual event
& 5-10\% of full program \\

\hline

Q2: PREVENT
& $\xi > 0.6$ $\delta \leq 0.6$
& Exponential
& Monthly + quarterly + targeted subsidy
& 30-50\% of full program \\

\hline

Q3: DEFEND
& $\xi \leq 0.6$ $\delta > 0.6$
& Exponential
& Weekly touchpoint + peer support + reinforcement
& 50-70\% of full program \\

\hline

Q4: DECIDE
& $\xi \leq 0.6$ $\delta \leq 0.6$
& Hyperbolic
& Intensive (weekly) + structural change OR abandon
& 80-100\% of full program OR 0\% \\

\hline
\end{tabular}
\end{center}

%=============================================================================
% SECTION 7: KEY RESULTS
%=============================================================================

\section{Key Results}
\label{sec:bb-results}

This appendix establishes the following key results for operationalizing equilibrium dynamics:

\begin{enumerate}
    \item \textbf{Result R1 (Stability Operationalization):}
    Equilibrium stability $\xi \in [0,1]$ can be reliably estimated from three observable
    components: Network Effects ($N_{\text{net}}$), Institutional Support ($I_{\text{struct}}$),
    and Path Dependence ($P_{\text{dep}}$) with interaction term.

    \item \textbf{Result R2 (Critical Threshold Formula):}
    The critical threshold $\kappa_{\text{crit}}$ is determined by five mechanisms (M1-M5)
    with opposing effects:
    \begin{equation}
    \kappa_{\text{crit}} = \kappa_0 - (w_1 + w_3 + w_4 + w_5) + w_2 + \epsilon_{\text{domain}}
    \end{equation}

    \item \textbf{Result R3 (Four-Quadrant Classification):}
    The $(\xi, \delta)$ space partitions into four distinct regimes (Q1-Q4) with
    qualitatively different $\lambda(\kappa)$ dynamics and intervention requirements.

    \item \textbf{Result R4 (Regime-Specific Decay):}
    \begin{itemize}[nosep]
        \item Q1 (Stable + Near): Polynomial decay --- low-intensity maintenance sufficient
        \item Q2/Q3 (Mixed): Exponential decay --- proactive intervention required
        \item Q4 (Labile + Far): Hyperbolic singularity --- immediate decision required
    \end{itemize}

    \item \textbf{Result R5 (Cost-Proportionality):}
    Intervention costs scale predictably: Q1 (5-10\%), Q2 (30-50\%), Q3 (50-70\%), Q4 (80-100\% or 0\%)
    of full program cost, enabling budget allocation under uncertainty.
\end{enumerate}

%=============================================================================
% SECTION 8: WORKED EXAMPLES
%=============================================================================

\section{Worked Examples: Four Domains}

\subsection{Example 1: Solar Adoption in Urban Neighborhood}

\textbf{Context:} City district with 30\% solar adoption rate; good policy support;
visible rooftop installations.

\textbf{Step 1 – Measure $\xi$:}
\begin{align}
Q_1 &= 5 \quad \text{(30\% adoption, high visibility)} \\
Q_2 &= 4 \quad \text{(moderate cascades; when one neighbor installs, others ask)} \\
Q_3 &= 5 \quad \text{(strong subsidies, grid connection, public support)} \\
Q_4 &= 4 \quad \text{(adoption behavior is 2-5 years old, fairly established)} \\
\xi &= 0.85 \quad \text{(Highly Stable)}
\end{align}

\textbf{Step 2 – Measure Mechanisms:}
\begin{align}
M_1 &= 5 \quad (\text{Very visible}) \\
M_2 &= 1 \quad (\text{One-time investment, not daily habit}) \\
M_3 &= 5 \quad (\text{Bill impact visible monthly}) \\
M_4 &= 3 \quad (\text{Somewhat identity: "environmentalist" or "savvy investor"}) \\
M_5 &= 5 \quad (\text{Full infrastructure support}) \\
\kappa_{\text{crit}} &= 0.30 - 0.15 + 0.15 - 0.15 - 0.15 - 0.15 = 0.23
\end{align}

\textbf{Step 3 – Classify Context:}
\begin{align}
\delta &= 0.30 / 0.50 = 0.60 \quad \text{(near optimum, maybe slightly below)} \\
\text{Quadrant} &\approx Q1 \text{ (MAINTAIN)}
\end{align}

\textbf{Step 4 – Intervention Strategy:}
\begin{itemize}
    \item Form: Polynomial decay
    \item Intervention: Monthly email updates, annual neighborhood solar festival
    \item Cost: ~\$50/household/month for maintenance
    \item Timeline: Indefinite, low urgency
    \item Message: "This is stable. Small investments in community engagement keep adoption momentum."
\end{itemize}

\subsection{Example 2: COVID Vaccination in Polarized Community}

\textbf{Context:} Community with 55\% vaccination rate; strong pro/anti-vax polarization;
vaccination certificates are controversial; no mandate.

\textbf{Step 1 – Measure $\xi$:}
\begin{align}
Q_1 &= 2 \quad \text{(55\% vaccination, but highly polarized; not reinforcing)} \\
Q_2 &= 2 \quad \text{(cascades in BOTH directions: pro-vax → more, anti-vax → more)} \\
Q_3 &= 2 \quad \text{(weak policy, no infrastructure, controversial)} \\
Q_4 &= 2 \quad \text{(vaccination is one-time event; identity is the ongoing feature)} \\
\xi &= 0.35 \quad \text{(Metastable – could go either way)}
\end{align}

\textbf{Step 2 – Measure Mechanisms:}
\begin{align}
M_1 &= 5 \quad \text{(Highly visible; core political/moral signal)} \\
M_2 &= 1 \quad \text{(One-time decision)} \\
M_3 &= 1 \quad \text{(No visible feedback; benefits are abstract/"invisible")} \\
M_4 &= 5 \quad \text{(CORE identity: "vaccinated person" vs "vaccine-hesitant person")} \\
M_5 &= 2 \quad \text{(Weak institutional support in this community)} \\
\kappa_{\text{crit}} &= 0.30 - 0.15 + 0.15 - 0.15 - 0.15 - 0.05 = 0.05 \quad \text{(VERY LOW!)}
\end{align}

\textbf{Step 3 – Classify Context:}
\begin{align}
\delta &= 0.55 / 0.85 = 0.65 \quad \text{(reasonably close to optimum)} \\
\text{Quadrant} &= Q2 \text{ or Q3 \quad (mixed; depends on how you weight $\xi$ vs $\delta$)}
\end{align}

\textbf{Step 4 – Intervention Strategy:}
\begin{itemize}
    \item Form: Hyperbolic (bifurcation at 0.05 is VERY early!)
    \item Interpretation: System is metastable. Could tip pro-vax OR anti-vax.
    \item Intervention:
    \begin{itemize}
        \item HIGH PRIORITY: Reinforce pro-vax identity (e.g., "vaccinated community leader" testimonials)
        \item HIGH PRIORITY: Counter anti-vax cascades (misinformation campaigns)
        \item Avoid: Mandate (polarizes further)
    \end{itemize}
    \item Cost: Medium-High
    \item Timeline: URGENT; tipping point could occur in 1-3 months
    \item Message: "This is critical. Identity-based reinforcement is essential; cascades can go either direction very quickly."
\end{itemize}

\subsection{Example 3: Smoking Cessation Program}

\textbf{Context:} Individual who quit smoking 3 months ago; has support group;
no strong institutional barriers to relapse (cigarettes freely available).

\textbf{Step 1 – Measure $\xi$:}
\begin{align}
Q_1 &= 1 \quad \text{(No peer influence; smoking is declining, not cool)} \\
Q_2 &= 1 \quad \text{(No cascades; quitting is individual)} \\
Q_3 &= 3 \quad \text{(Partial support: support group present, but no policy)} \\
Q_4 &= 2 \quad \text{(Quitting is recent; old "smoker identity" still strong)} \\
\xi &= 0.40 \quad \text{(Labile – fragile)}
\end{align}

\textbf{Step 2 – Measure Mechanisms:}
\begin{align}
M_1 &= 1 \quad \text{(Smoking is invisible socially)} \\
M_2 &= 5 \quad \text{(20+ years of daily smoking; DEEP habit)} \\
M_3 &= 2 \quad \text{(Benefits are delayed and abstract for ex-smoker)} \\
M_4 &= 4 \quad \text{(Identity is shifting: "smoker" → "person who quit" but still fragile)} \\
M_5 &= 1 \quad \text{(No institutional barriers to relapse; cigarettes are available)} \\
\kappa_{\text{crit}} &= 0.30 - 0.00 + 0.15 - 0.05 - 0.06 - 0.00 = 0.34
\end{align}

\textbf{Step 3 – Classify Context:}
\begin{align}
\delta &= 3\text{ months} / 6\text{ months goal} = 0.50 \quad \text{(midway to optimum)} \\
\text{Quadrant} &= Q3 \text{ (DEFEND)}
\end{align}

\textbf{Step 4 – Intervention Strategy:}
\begin{itemize}
    \item Form: Exponential (with potential bifurcation at $\kappa \approx 0.34$)
    \item Critical Phase: Month 2-4 of cessation is HIGH-RISK (extinction burst)
    \item Intervention:
    \begin{itemize}
        \item Weekly check-ins with counselor (month 1-6)
        \item Support group attendance (weekly)
        \item Cognitive reframing: reinforce "non-smoker identity"
        \item Relapse triggers: identify high-risk situations; plan coping strategies
    \end{itemize}
    \item Cost: High (weekly professional engagement)
    \item Timeline: 6 months intensive support; after that, monthly maintenance
    \item Message: "This is fragile. Deep habits have sharp extinction bursts. Support must be WEEKLY through month 4."
\end{itemize}

\subsection{Example 4: Workplace Wellness (Exercise Program)}

\textbf{Context:} Corporate fitness program; 40\% participation; gym access provided;
mixed enthusiasm; fall/winter attendance drops.

\textbf{Step 1 – Measure $\xi$:}
\begin{align}
Q_1 &= 3 \quad \text{(40\% participation; moderate peer influence)} \\
Q_2 &= 3 \quad \text{(Some cascades: gym buddy effects)} \\
Q_3 &= 4 \quad \text{(Strong: gym access free, time off allowed, company challenges)} \\
Q_4 &= 3 \quad \text{(Moderate: 1-2 years for active members)} \\
\xi &= 0.60 \quad \text{(Borderline Stable/Metastable)}
\end{align}

\textbf{Step 2 – Measure Mechanisms:}
\begin{align}
M_1 &= 3 \quad \text{(Moderate visibility: gym buddy networks)} \\
M_2 &= 3 \quad \text{(Moderate habit: weekly/monthly practice)} \\
M_3 &= 4 \quad \text{(Good feedback: visible fitness improvements in weeks)} \\
M_4 &= 2 \quad \text{(Weak identity: "I go to gym" but not core self-concept)} \\
M_5 &= 4 \quad \text{(Strong: gym free, time allocated, challenges organized)} \\
\kappa_{\text{crit}} &= 0.30 - 0.09 + 0.09 - 0.12 - 0.06 - 0.12 = 0.28
\end{align}

\textbf{Step 3 – Classify Context:}
\begin{align}
\delta &= 0.40 / 0.60 = 0.67 \quad \text{(reasonably close to optimum)} \\
\text{Quadrant} &= Q1 \text{ to Q2 \quad (close to maintaining, but risk in winter)}
\end{align}

\textbf{Step 4 – Intervention Strategy:}
\begin{itemize}
    \item Form: Polynomial to Exponential (seasonal shifts)
    \item Critical Period: October-December (attendance drops sharply)
    \item Intervention:
    \begin{itemize}
        \item Summer (May-Sept): Monthly challenges, gym buddy matching
        \item Fall (Oct-Nov): Increase frequency to WEEKLY challenges, indoor events
        \item Winter (Dec-Feb): Intensive engagement (2× per week), outdoor alternatives
        \item Spring (Mar-May): Ramp back down as external motivation rises
    \end{itemize}
    \item Cost: Moderate; varies by season (5-10\% in summer, 40-60\% in winter)
    \item Timeline: Annual cycle; critical decision point in September
    \item Message: "Stable baseline with seasonal drift risk. Proactive intensification in Q4 prevents January dropout."
\end{itemize}

%=============================================================================
% SECTION 8: INTEGRATION WITH CHAPTER 15
%=============================================================================

\section{Integration with Chapter 15: Behavioral Change Segments \& Journey Integration}

\subsection{Workflow: Diagnosis → Design → Deploy}

Chapter 15 integrates three dimensions: BCS (Segments), BCJ (Journey), AV (Willingness).
This Appendix BB provides the \textbf{diagnostic and strategic foundation}:

\begin{enumerate}
    \item \textbf{Phase 1 – Characterize Your Context (Appendix BB):}
    \begin{itemize}
        \item Use diagnostic questionnaire to estimate $(\xi, \kappa_{\text{crit}}, M_1,...,M_5)$
        \item Determine your quadrant (Q1-Q4)
        \item Select appropriate $\lambda(\kappa)$ form
    \end{itemize}

    \item \textbf{Phase 2 – Segment Your Population (Chapter 14 / Appendix BA):}
    \begin{itemize}
        \item Identify behavioral clusters (BCS types)
        \item Measure heterogeneity in willingness within segments
        \item Apply leverage profiles to find high-impact intervention points
    \end{itemize}

    \item \textbf{Phase 3 – Understand Journey Dynamics (Chapter 13 / Appendix AW):}
    \begin{itemize}
        \item Map where populations are in the Behavioral Change Journey
        \item Identify which phase (S1-S6) has highest dropout risk
        \item Match intervention timing to journey phases
    \end{itemize}

    \item \textbf{Phase 4 – Design Integrated Intervention (Chapter 15):}
    \begin{itemize}
        \item For each $(BCS_k, \varphi_j, W_{\text{eff}}(t))$ combination
        \item Recommend intervention mix based on your quadrant
        \item Allocate resources proportional to dropdown risk
    \end{itemize}

    \item \textbf{Phase 5 – Monitor \& Adapt (Ongoing):}
    \begin{itemize}
        \item Track $\kappa(t)$ monthly or quarterly
        \item If drift toward $\kappa_{\text{crit}}$ detected, increase intensity
        \item If approaching $\kappa_{\text{crit}}$, make critical decision
    \end{itemize}
\end{enumerate}

\subsection{The Feedback Loop: W11-W13 → Intervention → $\kappa$ Trajectory}

Recall W11-W13 (Appendix AV):

\begin{equation}
W_{\text{eff}}(t) = W_{\text{base}} \cdot (1 + \kappa_{AWX}(t)) \cdot (1 + \kappa_{KON}(t)) \cdot (1 + \kappa_{JNY}(t))
\end{equation}

\noindent
where:
\begin{itemize}
    \item $\kappa_{AWX}(t)$ = Awareness modulation (how salient is the behavior?)
    \item $\kappa_{KON}(t)$ = Context modulation (how conducive is the environment?)
    \item $\kappa_{JNY}(t)$ = Journey modulation (what phase are they in?)
\end{itemize}

\noindent
An intervention $I(t)$ works by \textbf{moving the modulation parameters}:

\begin{itemize}
    \item \textbf{Awareness Campaign} $\to$ increases $\kappa_{AWX}$
    \item \textbf{Environmental Design} $\to$ increases $\kappa_{KON}$
    \item \textbf{Temporal Structuring} $\to$ optimizes $\kappa_{JNY}$
\end{itemize}

\noindent
The combination determines $\kappa(t)$ trajectory and whether it crosses $\kappa_{\text{crit}}$.

%------------------------------------------------------------------------------
% HHH (METHOD-TOOLKIT) INTEGRATION - NEW JANUARY 2026
%------------------------------------------------------------------------------
\subsection{Integration with Appendix HHH (METHOD-TOOLKIT)}
\label{sec:bb-hhh-integration}

\begin{tcolorbox}[colback=orange!5!white, colframe=orange!75!black,
    title=\textbf{Integration: Equilibria $\leftrightarrow$ HHH (Intervention Toolkit)}]

\textbf{Why This Matters:}
Each intervention $I(t)$ in the equilibrium model corresponds to a HHH-compliant
intervention specification $\mathcal{I}$. The 20-field schema ensures that
interventions are systematically specified for equilibrium analysis.

\textbf{Mathematical Connection:}
\begin{equation}
I(t) = E(\mathcal{I}) \cdot \mathbb{1}[t \in [t_{\text{start}}, t_{\text{end}}]]
\end{equation}
where $E(\mathcal{I})$ = intervention effectiveness from HHH.

\textbf{Required HHH Fields for Equilibrium Analysis:}
\begin{center}
\renewcommand{\arraystretch}{1.2}
\begin{tabular}{@{}lll@{}}
\toprule
\textbf{HHH Field} & \textbf{Equilibrium Parameter} & \textbf{Function} \\
\midrule
F12 (Time Horizon) & $[t_{\text{start}}, t_{\text{end}}]$ & Intervention window \\
F14 (Path Function) & Equilibrium transition type & Door-opener/Amplifier/Stabilizer \\
F15 (Repetition) & Continuous vs. discrete $I(t)$ & Cyclic/one-time \\
F10 (Complementarity) & $\gamma_{ij}$ matrix & Portfolio interactions \\
\bottomrule
\end{tabular}
\end{center}

\textbf{Quadrant-Intervention Mapping (from HHH Types):}
\begin{center}
\renewcommand{\arraystretch}{1.2}
\begin{tabular}{@{}lll@{}}
\toprule
\textbf{Quadrant} & \textbf{Optimal HHH Types} & \textbf{Rationale} \\
\midrule
Q1 (High $\xi$, High $\kappa$) & Identity (7), Commitment (5) & Lock in stability \\
Q2 (High $\xi$, Low $\kappa$) & Structural (3), Incentive (4) & Reduce $\theta$ urgently \\
Q3 (Low $\xi$, High $\kappa$) & Normative (2), Feedback (6) & Build stability slowly \\
Q4 (Low $\xi$, Low $\kappa$) & Hybrid (9), intensive portfolios & Crisis mode \\
\bottomrule
\end{tabular}
\end{center}

\textbf{What Equilibria Provides to HHH:}
\begin{itemize}[nosep]
\item Critical threshold $\kappa_{\text{crit}}$ for intervention timing
\item Stability parameter $\xi$ for intensity calibration
\item Trajectory monitoring for adaptive intervention
\end{itemize}

\textbf{What HHH Provides to Equilibria:}
\begin{itemize}[nosep]
\item Standardized intervention specification (20 fields)
\item Complementarity constraints for portfolio design
\item Path function taxonomy for transition modeling
\end{itemize}

\textbf{Interface:} BB $I(t)$ specification $\leftrightarrow$ HHH $\mathcal{I}$ tuple

\textbf{Cross-Reference:} See Appendix HHH Section 3 (20-Field Schema) for complete specification.

\end{tcolorbox}

%=============================================================================
% SECTION 9: CRITICAL FOUNDATIONS
%=============================================================================

\section{Critical Foundations}

This Appendix BB rests on:

\begin{enumerate}
    \item \textbf{W11 (Dynamic Modulation):} Axiom that $\kappa$ is not static but dynamically
    modulated by AWX, KON, JNY in real-time

    \item \textbf{W12 (Boundary Conditions):} Axiom that Blockade ($A_{\text{mod}} = 0 \Rightarrow W = 0$)
    and Hyper-Awareness ($A_{\text{mod}} >> \theta$) create non-linear regimes

    \item \textbf{W13 (Bridge to Behavior):} Axiom that Willingness is input to probability
    function $P(i \to j, t) = g(W, \theta_{W,S}, ...)$, not deterministic

    \item \textbf{K-Axioms (Complementarity):} All modulation parameters $\kappa_{\text{AWX}}, \kappa_{\text{KON}}, \kappa_{\text{JNY}}$
    must respect complementarity constraints

    \item \textbf{Habit Extinction Literature:} Time constants for $\lambda(\kappa)$ decay are
    empirically grounded in behavioral psychology (Lally, Staddon, etc.)

    \item \textbf{Systems Dynamics / Bifurcation Theory:} Mathematical framework for
    understanding tipping points and phase transitions
\end{enumerate}

%=============================================================================
% SUMMARY
%=============================================================================

\section{Summary: From Diagnosis to Decision}

This appendix provides a complete toolkit for practitioners to operationalize behavioral change interventions:

\begin{itemize}
    \item \textbf{Diagnostic Stage:} Use the questionnaire (Section 6) to estimate equilibrium stability $\xi$ and critical threshold $\kappa_{\text{crit}}$ for your specific domain.

    \item \textbf{Classification Stage:} Map $(\xi, \delta)$ to one of four quadrants (Maintain, Prevent, Defend, Decide) determining intervention urgency.

    \item \textbf{Strategy Stage:} The quadrant determines both the $\lambda(\kappa)$ functional form (polynomial, exponential, or hyperbolic) and the optimal intervention mix.

    \item \textbf{Monitoring Stage:} Track $\kappa(t)$ over time; if approaching $\kappa_{\text{crit}}$, escalate intervention intensity or make critical go/no-go decision.

    \item \textbf{Integration Stage:} Combine diagnosis with Chapter 15's BCS × BCJ × AV 3D framework to design domain-specific interventions for each segment × journey-phase combination.
\end{itemize}

The framework bridges W11-W13 (dynamic willingness modulation) with practical intervention design, enabling evidence-based decision-making about when to maintain, prevent deterioration, defend status quo, or make critical investments in behavioral change.

%=============================================================================
% GLOSSARY & BIBLIOGRAPHY
%=============================================================================

\section{Glossary of Symbols}
\label{sec:bb-glossary}

See \textbf{Appendix G} (\textit{Glossary and Notation}) for comprehensive definitions.
Key symbols in this appendix:

\begin{itemize}[nosep]
    \item $\xi$ --- Equilibrium Stability parameter $\in [0,1]$
    \item $\kappa_{\text{crit}}$ --- Critical Threshold (Tipping Point)
    \item $\lambda(\kappa)$ --- Drift Rate Function
    \item $\delta$ --- Distance to Optimum ratio
    \item BCS, BCJ, AV --- Behavioral Change Segments, Journey, Willingness
    \item W11, W12, W13 --- Dynamic Modulation Axioms (from Appendix AV)
\end{itemize}

\textbf{Master Bibliography Reference:}

This appendix draws on literature indexed in:

\texttt{bibliography/bcm2\_master\_references.bib}

\nocite{bcm2_master}

%=============================================================================
% OPEN QUESTIONS & FUTURE WORK
%=============================================================================

\section{Open Questions for Future Research}

\begin{enumerate}
    \item \textbf{Domain-Specific Calibration:}
    The five mechanisms $(M_1,...,M_5)$ have assumed weights.
    How should these be adjusted for specific domains empirically?

    \item \textbf{Heterogeneity in $\kappa_{\text{crit}}$:}
    Do individuals within the same domain have different $\kappa_{\text{crit}}$ values
    based on psychology (risk tolerance, impulsivity, etc.)?

    \item \textbf{Intervention Efficacy Curves:}
    How does intervention type (Awareness, Context, Habit) map to different
    increases in $\kappa_{\text{AW}}, \kappa_{\text{KON}}, \kappa_{\text{JNY}}$?

    \item \textbf{Optimal Stopping Time:}
    When should practitioners decide to "abandon ship" rather than intensify
    intervention as they approach $\kappa_{\text{crit}}$?

    \item \textbf{Cost-Benefit Analysis:}
    Under what conditions is maintenance ($I_{\min}$) cost-justified versus
    letting the behavior collapse and respecting budget constraints?
\end{enumerate}

%=============================================================================
% REFERENCES
%=============================================================================

\section{References}

\begin{itemize}
    \item Appendix AV: Willingness Formalization (W11-W13)
    \item Appendix AW: Behavioral Change Journey (5-stage model)
    \item Chapter 14: Behavioral Change Segments (BCS operationalization)
    \item Chapter 15: Integration Framework (3D decision surface)
    \item Appendix BA: Formal Segmentation (axioms of segment emergence)
    \item Lally, P., et al. (2009). How are habits formed? European Journal of Social Psychology, 40(6).
    \item Staddon, J. E., et al. (2017). Habit and its neural substrate. Current Opinion in Behavioral Sciences.
\end{itemize}

\end{document}
